\chapter{Testing APIs with Postman}

Before wiring an endpoint up to React, test it in isolation. If a request
works correctly in Postman but fails once the frontend calls it, the bug is
in the frontend (or in how it calls the API); if it fails in Postman too,
the bug is in the backend. This separation saves significant debugging time
compared to only ever testing through the UI.

\begin{notebox}[NOTE]
For authenticated routes, enable Postman's cookie jar (it's on by default
in the desktop app) so the \texttt{HttpOnly} JWT cookie set by
\texttt{/login} is automatically stored and re-sent on subsequent requests
in the same collection --- exactly like a browser would.
\end{notebox}

\section{Register}

\textbf{POST} \texttt{http://localhost:5000/api/auth/register}\\
Headers: \texttt{Content-Type: application/json}

\begin{lstlisting}
// Body (raw, JSON)
{
  "name": "Ada Lovelace",
  "email": "ada@example.com",
  "password": "SecurePass123"
}
\end{lstlisting}

\begin{lstlisting}
// Expected response - 201 Created
{
  "success": true,
  "data": { "id": "6614...", "name": "Ada Lovelace", "email": "ada@example.com" }
}
\end{lstlisting}

\section{Login}

\textbf{POST} \texttt{http://localhost:5000/api/auth/login}\\
Headers: \texttt{Content-Type: application/json}

\begin{lstlisting}
// Body
{ "email": "ada@example.com", "password": "SecurePass123" }
\end{lstlisting}

Check the \textbf{Cookies} tab of the response in Postman --- you should see
the \texttt{token} cookie listed there. That confirms the cookie jar has it
stored for reuse on the next request.

\begin{lstlisting}
// Expected response - 200 OK
{ "success": true, "data": { "id": "6614...", "email": "ada@example.com" } }
\end{lstlisting}

\section{Get Profile (\texttt{/me})}

\textbf{GET} \texttt{http://localhost:5000/api/auth/me}\\
No body needed. No manual header needed either --- Postman automatically
resends the stored \texttt{token} cookie from the login step.

\begin{lstlisting}
// Expected response - 200 OK
{ "success": true, "data": { "id": "6614...", "name": "Ada Lovelace", "email": "ada@example.com" } }
\end{lstlisting}

If this returns \texttt{401 Unauthorized}, either the cookie jar did not
persist the cookie (check Postman settings) or the login step did not
actually set it (check the backend's \texttt{res.cookie} call).

\section{Create Task}

\textbf{POST} \texttt{http://localhost:5000/api/tasks}\\
Headers: \texttt{Content-Type: application/json}

\begin{lstlisting}
// Body
{ "title": "Write chapter 39", "description": "Postman testing", "status": "pending" }
\end{lstlisting}

\begin{lstlisting}
// Expected response - 201 Created
{
  "success": true,
  "data": { "_id": "6715...", "title": "Write chapter 39", "status": "pending", "owner": "6614..." }
}
\end{lstlisting}

\section{Get Tasks}

\textbf{GET} \texttt{http://localhost:5000/api/tasks}

\begin{lstlisting}
// Expected response - 200 OK
{
  "success": true,
  "data": [
    { "_id": "6715...", "title": "Write chapter 39", "status": "pending" }
  ]
}
\end{lstlisting}

\section{Update Task}

\textbf{PUT} \texttt{http://localhost:5000/api/tasks/6715...}\\
Headers: \texttt{Content-Type: application/json}

\begin{lstlisting}
// Body
{ "status": "done" }
\end{lstlisting}

\begin{lstlisting}
// Expected response - 200 OK
{ "success": true, "data": { "_id": "6715...", "title": "Write chapter 39", "status": "done" } }
\end{lstlisting}

\section{Delete Task}

\textbf{DELETE} \texttt{http://localhost:5000/api/tasks/6715...}\\
No body needed.

\begin{lstlisting}
// Expected response - 200 OK
{ "success": true, "message": "Task deleted" }
\end{lstlisting}

\section{What to Check on Every Request}

\begin{itemize}
\item[$\square$] Status code matches what the route is documented to return (200/201 for success, 4xx for expected client errors)
\item[$\square$] Response body has the expected shape (\texttt{success}, \texttt{data}/\texttt{message} keys)
\item[$\square$] \texttt{Content-Type: application/json} header is set on requests with a body
\item[$\square$] Cookie jar shows the \texttt{token} cookie after login, and it's being resent on later requests
\item[$\square$] Deleting/updating an ID that does not exist returns \texttt{404}, not a \texttt{500} crash
\end{itemize}

\begin{tipbox}[TIP]
Save each of these requests into a Postman collection with the base URL as
a collection variable (e.g.\ \texttt{\{\{baseUrl\}\}/tasks}). Switching from
local development to a deployed backend then only requires changing one
variable instead of every request.
\end{tipbox}

\whatwhyhow
{Every backend route is tested directly with Postman -- method, URL, headers, and body -- before any React code calls it, using Postman's cookie jar to persist the JWT across authenticated requests.}
{Testing through the UI conflates frontend and backend bugs; testing the API directly proves backend correctness first so any later frontend failure can be isolated to the frontend.}
{Send each request with the documented method/body, verify the status code and response shape, and confirm cookies persist across the login -> \texttt{/me} -> task CRUD sequence.}
